\documentclass{report}
\usepackage{amsmath,amssymb}
\setlength{\parindent}{0mm}
\setlength{\parskip}{1em}
\begin{document}
\begin{center}
\rule{6in}{1pt} \
{\large Kord Smith \\
{\bf Challenges and Approaches to Coupled Monte Carlo Neutronic/Fluid Flow Analysis of Commercial Nuclear Reactors}}

MIT \\ 77 Massachusetts Ave \\ Room 24-221 \\ Cambridge \\ MA 02139
\\
{\tt kord@mit.edu}\\
Kord Smith\\
Benoit Forget\\
Paul Romano\\
Bryan Herman\\
Andrew Siegel\end{center}



Despite the great advancements in Monte Carlo neutral particle transport
methods and their implementation on advanced computing platforms, Monte
Carlo methods have achieved decidedly little impact on either the design
or licensing of commercial light water reactors (LWRs). There are
numerous reasons that contribute to this lack of penetration:

1) Conventional Monte Carlo analysis requires 10�s to 100�s of billions
of neutron histories to resolve local spatial details in high dominance
ratio reactor cores
(e.g. source convergence is extremely difficult, even without multi-physics effects).

2) Traditional Monte Carlo statistical treatments lead to
non-conservative estimates of spatial parameter uncertainties because of
inter-cycle correlation effects in high dominance ratio cores (e.g.,
central limit theorem can not be applied directly).

3) Reactor core depletion requires tracking several hundred nuclides in
each radial ring of ~100,000 fuel pins at each of several hundred axial
levels (e.g., 100�s of billions of tallies are required).

4) Each fuel depletion region and each feedback region has unique physical properties
(e.g., nuclear data in each region must be treated as functions of local temperatures).

5) Accurate analysis requires dynamic coupling of Monte Carlo neutronics
to nuclear fuel/fluid flow models (e.g. coupled physics is crucial to
accuracy).

6) Steady-state core analysis cannot answer safety questions that require
dynamic multi-physics analysis (e.g., Monte Carlo methods must be
extended to the time domain for both rapid and long-duration transient
applications).

The authors will present a formal LWR cycle depletion benchmark that can
be used to test coupled neutronic/fluid flow models in pseudo
steady-state (reactor cycle depletion) applications. This benchmark will
be based on two cycles of operation of a commercial LWR, and measured
reactor data will be provided to assess computational model uncertainties
in core reactivity and 3D fission rate distributions. The inherent
challenges and difficulties in developing methods to solve this challenge
problem will be outlined and discussed in detail. In particular,
limitations of parallelization of Monte Carlo by replication and
challenges arising from massive storage requirements (for
temperature-dependent nuclear data and detailed spatial tallies) will be
discussed. A summary of recent algorithmic testing of alternative
parallelization strategies to overcome these limitations will be
presented.

The authors� experiences in applying one class of non-linear operators to
the acceleration of deterministic multi-physics methods will be
summarized, and computational results from recent extensions of these
methods to steady-state Monte Carlo analysis will be presented. Authors
will present results from recent parallel algorithm computational tests
and a proposed strategy for using non-linear operators to solve coupled
Monte Carlo neutronics and nuclear fuel/flow problems.

The authors will also present their proposed approach for extending these
non-linear operators to the development of transient Monte Carlo coupled
multi-physics reactor analysis methods. Key developments needed for
success will be outlined, and prospects for future success will be
discussed.


\end{document}
